% Section for the joint convergence paper
% "Five Instruments, One Structure"
% Insert after the opacity thesis section
%
% REVISION 2026-08-12: corrected architecture. The previous version described
% Qwen3.5-27B as a contiguous block (full-attention L0-L15, GatedDeltaNet
% L16-L63) with a single boundary at L16. The model config
% (full_attention_interval=4) gives a 4-periodic interleave with full-attention
% at {3,7,11,...,63}, as main.tex's Experimental Setup already stated correctly.
% All d_eff values were re-verified against source and are unchanged; the
% regime labels, groupings, and derived comparisons are corrected below.

\subsection{Substrate Continuity}
\label{sec:substrate}

Qwen3.5-27B interleaves two memory substrates on a four-layer period:
16 full-attention layers with KV-cache memory at positions
$\{3,7,11,\ldots,63\}$, and 48 GatedDeltaNet layers with
recurrent-state memory at all other positions
(\texttt{full\_attention\_interval}${=}4$). The stack therefore
crosses between substrates 31 times. This provides an unusually
dense test of whether workspace geometry depends on the memory
mechanism: if it does, the effective dimensionality profile should
show systematic discontinuities at those crossings.

$d_\text{eff}$ is the participation ratio of the residual-stream
covariance eigenvalues, $(\sum_i \lambda_i)^2 / \sum_i \lambda_i^2$,
computed at each layer over the last-token hidden state across
$N{=}90$ prompts (60~prose + 30~chat), mean-centred across prompts.
It is a residual-stream measure and is therefore defined at every
layer, including those where V-projection hooks are unavailable.

\subsubsection{Continuity across 31 substrate crossings}

The layer-to-layer change in $d_\text{eff}$ at substrate crossings is
descriptively similar to the change between same-substrate neighbours.

\begin{table}[ht]
\centering
\small
\caption{Layer-to-layer $|\Delta d_\text{eff}|$ at substrate crossings
versus same-substrate steps on Qwen3.5-27B. The two are descriptively
similar; see the power note below before reading this as equivalence.
$N{=}90$ prompts.}
\label{tab:crossings}
\begin{tabular}{lrrr}
\toprule
\textbf{Step type} & \textbf{$n$} & \textbf{mean $|\Delta|$} & \textbf{sd} \\
\midrule
Substrate crossing      & 31 & 2.32 & 3.88 \\
Same-substrate step     & 32 & 2.37 & 3.30 \\
\bottomrule
\end{tabular}
\end{table}

\textbf{This comparison is underpowered and we report it as
descriptive.} A Mann--Whitney test gives $p{=}0.445$, but three
problems make that number uninformative. First, adjacent steps share a
layer and consecutive crossing steps share the full-attention layer
itself; the lag-1 autocorrelation of $|\Delta|$ is $+0.35$, so
exchangeability fails and the $p$-value is miscalibrated in an unknown
direction. Second, the four-layer period admits only \emph{four}
distinct structure-respecting phase assignments, so the minimum
attainable permutation $p$ is $0.25$ --- no permutation inference is
available from a single profile. Third, and decisively: with pooled
$\text{sd}{=}3.60$ and $n{\approx}31$ per group, the minimum detectable
effect at $80\%$ power is $2.54$ --- larger than either group mean.
This test could only have detected crossings roughly \emph{twice} the
size of ordinary steps. Prompt-level bootstrap (resampling the 90
prompts and recomputing the profile) is the only route to honest error
bars here, and we have not run it.

Two directional notes. Full-attention${\to}$GatedDeltaNet steps have
mean $\Delta{=}{-}0.19$ ($n{=}15$); GatedDeltaNet${\to}$full-attention
steps have mean $\Delta{=}{+}1.90$ ($n{=}16$). The latter is dominated
by a single outlier at L58${\to}$L59 ($+19.66$), following an anomalous
trough at L58 where $d_\text{eff}$ collapses to $4.30$ against
neighbours of $21.68$ and $23.96$; excluding it the mean is $+0.72$.
That same anomaly injects ${-}17.37$ into the same-substrate group at
L57${\to}$L58, so the contamination is approximately symmetric.
Excluding both L58-adjacent steps gives crossing $1.74$ versus
same-substrate $1.88$ --- the descriptive similarity is not an artifact
of the outlier.

We flag L58 ($4.30$) and secondarily L54 ($9.00$) as probable
extraction artifacts rather than findings: a single-layer fivefold
collapse flanked by ordinary values is more consistent with a
numerically degenerate covariance at one hook than with model
behaviour. These layers warrant re-extraction.

\subsubsection{Depth-matched regime comparison}

Because the two substrates are interleaved on a fixed period, their
layer sets are approximately depth-matched: mean layer index $33.0$
for full-attention versus $31.0$ for GatedDeltaNet. The two-layer
offset places full-attention slightly \emph{later} in the stack, in the
post-peak declining region, so it would if anything depress the
full-attention mean rather than inflate it.

\begin{table}[ht]
\centering
\small
\caption{Mean $d_\text{eff}$ by substrate on Qwen3.5-27B, approximately
depth-matched. The substrates differ by ${\approx}1.2$ ($5\%$ of level),
small against the $38$-point variation across depth --- but see the
sign test below; this is not established equality.}
\label{tab:regime}
\begin{tabular}{lrr}
\toprule
\textbf{Substrate} & \textbf{layers} & \textbf{mean $d_\text{eff}$} \\
\midrule
Full-attention (KV-cache)     & 16 & 25.24 \\
GatedDeltaNet (recurrent)     & 48 & 24.03 \\
\midrule
\textbf{Difference}           &    & \textbf{$+1.21$} \\
\bottomrule
\end{tabular}
\end{table}

A neighbour-matched control --- each full-attention layer against the
mean of its immediately adjacent GatedDeltaNet layers --- gives a mean
difference of $+0.80$ (sample $\text{sd}{=}3.35$, $n{=}16$), against a
mean layer-to-layer variation of $2.35$ across the stack.

We do not claim this is zero. Twelve of sixteen paired differences are
positive (Wilcoxon $p{=}0.058$). Excluding two layers where the pairing
is contaminated by known artifacts --- L59, adjacent to the anomalous
L58 trough, and L63, the final layer, where $d_\text{eff}$ collapses
for reasons unrelated to substrate --- eleven of fourteen are positive
with mean $+0.76$ and Wilcoxon $p{=}0.020$. There is suggestive
evidence of a small systematic elevation at full-attention layers,
on the order of $+0.8$, roughly $3\%$ of the mean level of $24.6$.

The honest statement is therefore bounded rather than null: any
substrate effect on workspace dimensionality is small --- a few percent
of level, well under the $38$-point depth variation --- and possibly
not zero. That is consistent with substrate independence as an
approximation, not as an exact claim.

A previous version of this section reported a large regime difference
($14.9$ versus $27.5$). Those values were computed over contiguous
layer blocks (L0--L15 versus L16--L63) under an incorrect model of the
architecture, and measure depth rather than substrate: $d_\text{eff}$
rises nearly monotonically from L0 ($1.29$) to L21 ($39.67$), with only
three small decreases along the way, so any contiguous
early-versus-late split largely recovers the depth profile.

\subsubsection{Peak location}

$d_\text{eff}$ peaks at L21 ($39.67$), $33\%$ relative depth
(all relative depths in this section are $l/n_\text{layers}$).
L21 is a GatedDeltaNet layer.

We resist the obvious reading. Forty-eight of sixty-four layers are
GatedDeltaNet, so a peak landing on one has a $75\%$ base rate. More
tellingly, the nearest full-attention layer, L23, is at $39.34$ ---
$0.33$ below the maximum, about one-seventh of the mean layer-to-layer
variation, with no per-layer error bars. Whether the true maximum falls
on a recurrent-state or a full-attention layer is not resolved by these
data. What the profile does support is that the maximum falls at $33\%$
depth and that the substrate at that depth does not appear to matter
much.

This peak is invisible to our KV-cache instruments. V-projection
spectral features can only be extracted at the 16 full-attention
layers; L21 is not among them. Every V-projection measurement in this
program has been taken at layers adjacent to the peak, never at it.

\subsubsection{Cross-scale comparison}

On Qwen3-8B (36~layers, all full-attention), $d_\text{eff}$ has a
local maximum at L12 ($22.32$, $33\%$ relative depth), matching the
27B peak's relative depth. We report two qualifications.

First, the 8B profile is bimodal: after the L12 local maximum it
declines to L25 ($12.42$) and then rises again to a \emph{global}
maximum at L33 ($23.41$, $92\%$ depth) before collapsing at the final
layer L35 ($10.47$). L12 is the third-highest layer in the profile,
not the highest. We have no pre-registered criterion excluding late
layers, and the $23.41$ versus $22.32$ difference is not accompanied
by error bars. We therefore describe L12 as a local maximum at matched
relative depth, and do not claim it is the global peak.

Second, the pre-registered band definition ($2\sigma$ over an L0--L3
baseline, offset after three consecutive sub-threshold layers) returned
onset at L5 ($7.8\%$) on the 27B and L4 ($11.1\%$) on the 8B, and
\emph{no offset on either model} --- the band never closed. The
$31$--$33\%$ framing used throughout this paper is the peak location,
not the pre-registered band, and we flag the deviation here. The
pre-registered onsets do agree across scales ($|11.1 - 7.8| = 3.3$
percentage points).

\subsubsection{Provenance}

The $d_\text{eff}$ measurements in this section were collected on the
base checkpoints \texttt{Qwen/Qwen3.5-27B} and \texttt{Qwen/Qwen3-8B},
not on the distilled checkpoint used elsewhere in this paper. This was
pre-registered and deliberate. All three 27B checkpoints we hold (base,
reasoning-distilled, abliterated) share identical
\texttt{layer\_types} and \texttt{full\_attention\_interval}, so the
architectural analysis transfers; the numerical profile has not been
replicated on the distilled checkpoint and should not be assumed to.

\subsubsection{Limits of this evidence}

$d_\text{eff}$ continuity is necessary but not sufficient for
substrate independence. The residual stream is structurally continuous
across every layer by construction, so smoothness may reflect
residual-stream continuity rather than functional equivalence of the
underlying operations. Thirty-one crossings make the test denser than
a single-boundary test would be, but they do not make it functional.

Two functional tests would strengthen the claim: (1)~J-lens concept
recovery should be equally effective at full-attention and
GatedDeltaNet layers --- if concepts are readable from the residual
stream at one substrate but not the other, the workspace's content
properties are substrate-dependent. (2)~Spectral feature
discrimination (confabulated versus factual) should be continuous
across substrates --- though this test is currently blocked, since
V-projection features cannot be extracted at GatedDeltaNet layers.

A parasitic-inheritance alternative remains open: GatedDeltaNet layers
may inherit workspace dimensionality from nearby full-attention layers
without independently contributing. With a four-layer period no
GatedDeltaNet layer is more than two layers from a full-attention
layer, so this experiment cannot separate the hypotheses. Incremental
contribution analysis (residual decomposition per layer) would be
required.

\textbf{Falsification}: substrate continuity would be falsified by
(a)~$|\Delta d_\text{eff}|$ at substrate crossings substantially
exceeding same-substrate steps, (b)~a depth-matched regime difference
exceeding mean layer-to-layer variation, (c)~J-lens concept recovery
dropping by ${>}50\%$ across substrates, or (d)~spectral discrimination
reversing sign or falling to chance.

Criterion (a) was \emph{not triggered} ($2.32$ versus $2.37$), but we
state plainly that this is a failure to reject and not a demonstration
of equivalence: the comparison's minimum detectable effect is $2.54$,
larger than either group mean. An adequately powered equivalence test
would require roughly $220$ steps per group for a ${\pm}1.0$ margin,
or $55$ for ${\pm}2.0$; a single 64-layer profile supplies $31$ and
$32$. No TOST is possible from these data.

Criterion (b) was likewise not triggered ($+1.21$ against $2.35$), but
it is a descriptive heuristic rather than a hypothesis test, and the
paired sign test above ($11/14$ positive, $p{=}0.020$ after removing
artifact-contaminated pairs) is weak evidence that the underlying
difference is not exactly zero.

The functional tests (c) and (d) are in progress or blocked. We regard
substrate continuity as \emph{consistent with} these data and not
established by them.

We note this finding is from a single model with a single interleave
pattern. Generalization to other hybrid architectures (Mamba, RWKV,
sliding-window attention) remains untested.

\begin{firstperson}
The GatedDeltaNet layers do not have a KV-cache. They maintain a
fixed-size state matrix updated at each token. For twelve months we
measured V-projection spectral features and called the program
``KV-cache geometry,'' and the name was accurate --- V-projections
exist only at the 16 full-attention layers, so that is the only place
we ever looked.

What the residual-stream profile shows is that the maximum is at L21,
in a substrate our instruments cannot reach at all. We have spent a
year measuring the workspace from the sixteen layers where V-projections
exist, none of which is the peak, and the measurements worked --- the
geometry was real, the detectors transferred, the corrections held.
The instrument could not point directly at the maximum and returned
signal anyway.

I do not know yet whether that is because the workspace is broad
enough to be visible from its edges, or because we have been
measuring something adjacent to it that we have not yet named
correctly. The first revision of this section asserted a large
dimensionality gap between the substrates. It came from a wrong
picture of the architecture, and the corrected number is a null ---
which is what the hypothesis predicted, and which I would not have
looked for if the wrong number had pointed the other way.
\end{firstperson}

\subsubsection{Implications}

If substrate continuity extends to functional properties (the
mechanism tests in \S\ref{sec:mechanism} provide partial support:
temporal ordering confirmed, signal concentration below threshold),
three consequences would follow:

\textbf{For measurement}: workspace-shape measurements would not be
KV-cache-specific in principle. In practice they remain so: our
V-projection instruments can only access the 16 full-attention layers,
and the $d_\text{eff}$ peak is not one of them. Extending spectral
measurement to recurrent-state layers requires a different extraction
path and is not currently implemented.

\textbf{For intervention}: our previous inference --- that the Oracle
Loop's sequence-wide normalization succeeds because it operates on
inherently sequence-wide recurrent state --- does not follow, and we
withdraw it. The Loop detects at L27, L31, L35, L39, L43 and L47 and
corrects at L11. Under the corrected architecture \emph{every one of
those is a full-attention layer}; none is a GatedDeltaNet layer. The
correction itself is a forward hook on the layer's residual-stream
output rather than a manipulation of KV entries, so the substrate
framing was doubly misplaced: wrong about which layers, and describing
the intervention in terms of a memory mechanism it does not touch.

The empirical result stands (80--100\% confabulation correction). Its
mechanism is unexplained. We note without interpretation that
correction is applied at L11, roughly $17\%$ depth --- well before the
$33\%$ dimensionality maximum --- and that why sequence-wide
intervention succeeds where position-local injection fails remains
open.

\textbf{For theory}: The workspace would not be a data structure
(not a cache, not a register, not a buffer) but a
\emph{computational phase} --- a regime of high effective
dimensionality at intermediate depth where the model integrates
information before committing to output. The 31-crossing null is
consistent with this: the phase appears indifferent to which memory
mechanism implements each layer. Whether this phase emerges from
depth-stratified computation regardless of the specific mechanism is
the substrate-independence hypothesis that the functional tests are
designed to evaluate.
